% Generated draft: requires mathematical and source review.
\documentclass[11pt]{article}
\usepackage[margin=25mm]{geometry}
\usepackage{fontspec}
\setmainfont{Latin Modern Roman}
\usepackage{amsmath,amssymb,mathtools,unicode-math}
\setmathfont{Latin Modern Math}
\usepackage{enumitem,needspace,xurl,hyperref}
\setlength{\parindent}{0pt}
\setlength{\parskip}{5pt}
\setlength{\emergencystretch}{3em}
\newcommand{\N}{\mathbb{N}}
\newcommand{\Z}{\mathbb{Z}}
\newcommand{\Q}{\mathbb{Q}}
\newcommand{\R}{\mathbb{R}}
\newcommand{\C}{\mathbb{C}}
\newcommand{\F}{\mathbb{F}}
\newcommand{\A}{\mathbb{A}}
\newcommand{\PP}{\mathbb{P}}
\newcommand{\E}{\mathbb{E}}
\newcommand{\eps}{\varepsilon}
\newcommand{\dd}{\,\mathrm{d}}
\begin{document}
\section*{UnsolvedMath ID: 30001873}
\textbf{Source problem code:} catalogue entry
\subsection{Definitions and mathematical statement}
Under the usual scalar real-valued interpretation, let \((Y_n)_{n\in\Z}\) be a Gaussian process on a probability space \((\Omega,\mathcal F,\mathbb P)\): every finite vector of its coordinates has a multivariate Gaussian distribution. The supplied statement does not explicitly specify the scalar field, centering, or nondegeneracy. Stationarity means that its joint distributions are invariant under integer shifts. Ergodicity means that every shift-invariant event for its law on the sequence space has probability zero or one.

A prediction rule is interpreted as a sequence \(g=(g_n)_{n\geq1}\) of deterministic Borel measurable functions \(g_n:\R^{n-1}\to\R\), where \(\R^0\) is a singleton. Write \(\mathcal F_{n-1}^{+}=\sigma(Y_1,\ldots,Y_{n-1})\), with \(\mathcal F_0^{+}\) the trivial sigma-algebra. Gaussian coordinates are integrable, so the conditional expectations below exist.

The question, retaining the supplied quantifier order, is: for every stationary ergodic Gaussian process \((Y_n)_{n\in\Z}\), does there exist a prediction rule \(g\) such that
\[
\mathbb P\!\left(\lim_{n\to\infty}\left|\mathbb E[Y_n\mid\mathcal F_{n-1}^{+}]-g_n(Y_1,\ldots,Y_{n-1})\right|=0\right)=1?
\]
Although the process is indexed by all integers, the predictor and its conditional-expectation target use only the finite observed history starting at index \(1\).

\emph{Unresolved qualification.} The written order is \(\forall (Y_n)\,\exists g\), which permits the rule to depend on the process law. The record does not specify whether such dependence is allowed or whether the intended question instead asks for one rule, independent of the unknown law, that works for every process. That latter requirement would have the different order \(\exists g\,\forall (Y_n)\). No choice between these targets is made here. The real-valued and measurable-rule conventions above are explicit interpretations requiring review against the source.
\subsection{Short English statement}
For each stationary ergodic Gaussian process, the supplied question asks whether predictions based on \(Y_1,\ldots,Y_{n-1}\) can approach the corresponding conditional mean almost surely as \(n\to\infty\). It leaves unclear whether the prediction rule may depend on the process law or must be one universal rule. The supplied open status is dated curation, not an independently verified conclusion.
\subsection{Sources}
\begin{itemize}
\item Local UnsolvedMath registry, supplied complete record: permanent ID \(30001873\), exact source problem code \texttt{catalogue entry}, title ``Strong Prediction for Stationary Ergodic Gaussian Processes,'' collection ``Oberwolfach Reports Open Problems.'' No individual proposer is supplied; the proposed year is recorded as \(2011\).
\item Supplied source citation: \emph{Mini-Workshop: Mathematics of Machine Learning (2012)}, Oberwolfach Reports, pp.\ 2365--2396, DOI \texttt{10.4171/owr/2011/42}. Original OWR report: \url{https://doi.org/10.4171/owr/2011/42}. The supplied background describes these reports as documenting mathematical workshops and their open-problem discussions.
\item Supplied literature reference, labelled ``strongly consistent learning of the best predictor'': \url{https://pure.royalholloway.ac.uk/files/23266714/gysafinal.pdf}. Authors and further bibliographic details are not supplied.
\item Literature assessment dated exactly \texttt{2026-08-21}: the record labels the problem open and states that strong almost-sure prediction for every stationary ergodic Gaussian process is explicitly left open, while only Cesaro consistency and results under extra assumptions are known. This is supplied dated curation derived from the OWR curation, not independently verified here; its status and cited claims require independent verification.
\item The supplied statement status is \texttt{corrected\_verified}. Its verification note reports comparison with the official Oberwolfach report and normalization of notation, context, and question form without changing the source-level claim. This is a supplied verification assertion, not a source consultation performed for this draft. The citation's workshop year \(2012\), the proposed year \(2011\), and the DOI containing \(2011\) are retained without reconciliation.
\end{itemize}
\end{document}
